\documentclass[12pt,a4paper,oneside]{article}
\usepackage[brazil]{babel}
\usepackage[utf8]{inputenc}
\usepackage{olymp}

\setcounter{problem}{1}

\begin{document}

\begin{problem}{Sapiens}{sapiens}{2}
 
Faça um programa para ajudar o SAPIENS a determinar a situação de um estudante (aprovado, prova final ou reprovado) na disciplina INF110 ao final do semestre. Para isso, seu programa deve analisar a nota final do estudante ($NF$), seu número de faltas em aulas teóricas ($FT$) e seu número de faltas em aulas práticas ($FP$). As regras são as seguintes:

\begin{itemize}
    \item Todo estudante que tenha mais de 15 faltas em aulas teóricas está reprovado;
    \item Todo estudante que tenha mais de 7 faltas em aulas práticas está reprovado;
    \item Caso não tenha sido reprovado por faltas, sua situaçao será determinada pela sua nota final:
    \begin{itemize}
        \item $NF < 40 \Rightarrow$ Reprovado
        \item $40 \leq NF < 60 \Rightarrow$ Prova final
        \item $NF \geq 60 \Rightarrow$ Aprovado
    \end{itemize}
\end{itemize}

\Notes
%
Cada falta corresponde a uma aula de 1 hora de duração, ou seja, um dia de aula perdido implica em 2 faltas no SAPIENS! Aqui já será informado o total de faltas.


\InputFile

A entrada contém três valores inteiros, $NF$, $FT$ e $FP$, conforme descrito acima. Restrições: $0 \leq NF \leq 100$, $0 \leq FT \leq 60$, $0 \leq FP \leq 30$.

%\Notes
%Neste problema, é garantido que sempre haverá duas raízes reais.


\OutputFile

Seu programa deve gerar apenas uma linha na saída, contendo uma das seguintes palavras, dependendo da situação do estudante: ``\texttt{Aprovado}'', ``\texttt{Final}'' ou ``\texttt{Reprovado}''.

\Example

\begin{example}
\exmp{
100 0 8
}{
Reprovado
}%
\end{example}


\begin{example}
\exmp{
60 14 6
}{
Aprovado
}%
\end{example}

\begin{example}
\exmp{
59 0 0
}{
Final
}%
\end{example}

\begin{example}
\exmp{
53 8 8
}{
Reprovado
}%
\end{example}


%Comentários sobre o exemplo, se necessário.

\end{problem}

\end{document}
